\begin{pastexam}[label={px:ntu-109-dsaa-11}]{臺大資工所 109 學年度 第 11 題}
（問答）

(5 points) Given an array of $n$ integers $A=(a_1,a_2,\ldots,a_n)$ ($n>0$), you are asked to find
the maximum and minimum elements with minimum comparisons. A simple solution is to iteratively
compare each element with current max and min. It takes $2n$ comparisons in total.
\begin{enumerate}[label=(\alph*)]
\item (3 points) Devise an algorithm that requires less than $1.67n$ comparisons.
\item (2 points) How many comparisons does your algorithm take? Show your derivation.
\end{enumerate}
\end{pastexam}

\begin{pastsoln}
答案：兩兩一組，先組內比，大的去和目前的最大值比、小的去和目前的最小值比；總共
$\lceil 3n/2 \rceil - 2$ 次，小於 $1.5n < 1.67n$。

\textbf{(a) 演算法}（資結書第~16.1.2~節）。 簡單做法每個元素都要和最大值、最小值各比一次。 改進的想法：
兩個元素 $x$、$y$ 先互比一次，大的那個才有可能是最大值，小的那個才有可能是最小值，所以大的只需要和
目前的最大值比、小的只需要和目前的最小值比。 每兩個元素花 $3$ 次，而不是 $4$ 次。

\begin{pseudo}{同時找最大和最小（兩兩一組）}
\Function{Max-Min}{$A, n$}
    \If{$n$ 是奇數}
        \State $\mathit{max} \gets a_1$；$\mathit{min} \gets a_1$；$i \gets 2$
    \Else
        \If{$a_1 > a_2$} \Comment{第 $1$ 次比較}
            \State $\mathit{max} \gets a_1$；$\mathit{min} \gets a_2$
        \Else
            \State $\mathit{max} \gets a_2$；$\mathit{min} \gets a_1$
        \EndIf
        \State $i \gets 3$
    \EndIf
    \While{$i < n$} \Comment{每次處理 $a_i$、$a_{i+1}$ 一對}
        \If{$a_i > a_{i+1}$}
            \State $\mathit{big} \gets a_i$；$\mathit{small} \gets a_{i+1}$
        \Else
            \State $\mathit{big} \gets a_{i+1}$；$\mathit{small} \gets a_i$
        \EndIf
        \If{$\mathit{big} > \mathit{max}$}
            \State $\mathit{max} \gets \mathit{big}$
        \EndIf
        \If{$\mathit{small} < \mathit{min}$}
            \State $\mathit{min} \gets \mathit{small}$
        \EndIf
        \State $i \gets i + 2$
    \EndWhile
    \State \Return $(\mathit{max}, \mathit{min})$
\EndFunction
\end{pseudo}

正確性：迴圈每一輪開始時，$\mathit{max}$、$\mathit{min}$ 是 $a_1..a_{i-1}$ 的最大值和最小值。 一開始成立（奇數時只有
$a_1$；偶數時是 $a_1$、$a_2$）。 處理一對 $a_i$、$a_{i+1}$ 時，這一對的最大值是 $\mathit{big}$、最小值是 $\mathit{small}$，
$a_1..a_{i+1}$ 的最大值是 $\max(\mathit{max}, \mathit{big})$、最小值是 $\min(\mathit{min}, \mathit{small})$，第 15--18 行正好更新成這兩個。
迴圈結束時 $i = n + 1$，$\mathit{max}$、$\mathit{min}$ 是整個陣列的最大值和最小值。

\textbf{(b) 比較次數}。 迴圈裡每一對做 $3$ 次比較（第 11、15、17 行）。
\begin{itemize}
\item $n$ 是偶數：開頭 $1$ 次，剩下 $n - 2$ 個元素組成 $(n - 2)/2$ 對，共
\[
1 + 3 \cdot \frac{n - 2}{2} = \frac{3n}{2} - 2 .
\]
\item $n$ 是奇數：開頭 $0$ 次，剩下 $n - 1$ 個元素組成 $(n - 1)/2$ 對，共
\[
3 \cdot \frac{n - 1}{2} = \frac{3n - 3}{2} = \left\lceil \frac{3n}{2} \right\rceil - 2 .
\]
\end{itemize}
合起來是 $\lceil 3n/2 \rceil - 2$ 次（資結書定理~16.1），兩種情形都 $< 1.5n < 1.67n$。 例如
$n = 10$ 時是 $13$ 次，簡單做法要 $2(n - 1) = 18$ 次（題目說的 $2n$ 是粗略的算法）。

補充：$\lceil 3n/2 \rceil - 2$ 也是下界，任何用比較的演算法最壞都至少要這麼多次（資結書第~16.1.2~節的
對手論證），所以這個做法已經最好。
\end{pastsoln}
